\chapter{Concrete, Metals, and Fertilizer}


\section{Fertilizer}
One of the biggest problems humans face is: how can we get enough food to feed everyone?
In 1950, there were 2.5 billion people on the planet, and about 65\%
were \index{malnutrition}malnourished. In 2019, there were 7.7 billion people on the
planet, and only 15\% are malnourished. How did crop yields increase
so much? There were several factors: better crop varieties,
reliable irrigation, increased mechanization, and affordable fertilizers.\index{fertilizer}

When a plant grows, it takes molecules out of the soil and uses them
to build proteins. It primarily needs the elements nitrogen ($N$),
phosphorus ($P$), and potassium ($K$).\index{nitrogen} \index{phosphorus} \index{potassium}

When you buy a bag of fertilizer at the store, it typically has
three numbers on the front. For example, you might buy a bag of
``24-22-4''.  This means that 24\% of the mass of the bag is nitrogen. The other two numbers give phosphate (P$_2$O$_5$) and potash (K$_2$O), which are the forms that fertilizer labels use. The phosphorus and potassium in the bag are smaller amounts than these numbers suggest.

Potassium comes as potassium carbonate ($K_2CO_3$), potassium chloride
($KCl$), potassium sulfate ($K_2 SO_4$), and potassium nitrate
($KNO_3$). Any blend of these chemicals is known as ``potash''. Potash
is dug up out of mines. \index{potash}

Phosphorus is also mined, but is refined into phosphoric acid
($H_3PO_4$) before it is put into fertilizer.

Nitrogen is an especially interesting case for 2 reasons:
\begin{itemize}
\item Worldwide, farmers apply more nitrogen to their soil than potassium or phosphorous combined.
\item 78\% of the air we breathe is nitrogen in the form of $N_2$, but
  neither plants nor animals can utilize nitrogen in that form.
\end{itemize}

\subsection{The Nitrogen Cycle}

Converting the $N_2$ in the air into a form that a plant can use
is known as \newterm{nitrogen fixation}. For billions of years, there
were only two ways that nitrogen fixation occurred on earth:
\begin{itemize}
\item The energy from lightning causes $N_2$ and $H_2O$ to reconfigure as \index{ammonia}ammonia ($NH_3$) and nitrate ($NO_3$). This accounts for about 10\% of all naturally occurring nitrogen fixation.
\item Cyanobacteria are responsible for the rest. They convert $N_2$ into ammonia.
\end{itemize}\index{nitrogen cycle} \index{nitrogen fixation}

Let us say that you are eating soybeans. There is a cyanobacteria
called \newterm{rhizobia}\index{rhizobia} that has a \index{symbiosis}symbiotic relationship with
soybean plants. Rhizobia fixes nitrogen for the soybean plant. The
soybean plant performs \index{photosynthesis}photosynthesis, then gives sugars to the
rhizobia.

The proteins in the soybeans contain nitrogen from the rhizobia. When
you eat them, you use some of the nitrogen to build new proteins. You
probably do not use all the nitrogen, so your cells release ammonia into your blood.

Ammonia likes to react with things, so your liver combines the ammonia
with carbon dioxide to make urea ($CO(NH_2)_2$). Your kidneys take
the urea out of your blood and mix it with a bunch of water and salts
in your bladder.  When you urinate, the urea leaves your body.\index{urea}

If you urinate on the ground, the nearby plants can take the nitrogen out of
the urea.\index{urine}

When you die, the nitrogen in your proteins will return to the soil as
ammonia and nitrate.

For centuries, farms got their nitrogen from urine, feces, and rotting
organic material. There were two challenges with this:
\begin{itemize}
\item Human pathogens had to be kept away from human food.
\item There was simply not enough to support 7.7 billion people.
\end{itemize}

This meant we had to figure out how to do nitrogen fixation at an industrial
level.

\subsection{The Haber-Bosch Process}

During World War I, two German scientists, Fritz Haber and Carl Bosch,
figured out how to make ammonia from $N_2$ and $H_2$ using high
temperatures and pressures. This is how nearly all nitrogen fertilizer
is created today.\index{Haber-Bosch process}

Where do we get the $H_2$? From methane ($CH_4$) in \index{natural gas}natural gas. Today, 3-5\%
of the world's natural gas production is consumed in the Haber-Bosch
process.

The ammonia is converted into ammonium nitrate ($NH_4NO_3$) or urea
before it is shipped to farms.

\subsection{Other nutrients}

Healthy plants require several other elements that are sometimes
applied as fertilizer: calcium, magnesium, and sulfur.

Finally, tiny amounts of copper, iron, manganese, molybdenum, zinc, and
boron are sometimes needed.

%FIXME add photosynthesis process here?


\section{Concrete}
To make concrete, you mix cement with water and an \index{aggregate}aggregate (sand or
rock).  The cement is usually only about 10 to 15 percent of the
mixture. The cement reacts with the water, and the resulting solid
binds the aggregate together. In 2019, the world consumed 4.5 billion
tons of cement.\index{concrete} \index{cement}

Concrete is hard and durable. The mortar between the pyramids at Giza
is concrete --- it is now 5000 years old. Today, we use concrete to
build many structures including buildings, bridges, airport runways,
and dams. \index{mortar}Grout and mortar are related materials but different in that 
they are used as bonding or filler materials within construction.


There are many kinds of cement, but the most common is \index{Portland cement}Portland
cement. It is made by heating limestone (calcium carbonate) with clay
(for silicon) in a kiln. Two things come out of the kiln: Carbon
dioxide and a hard substance called ``\index{clinker}clinker''.  The clinker is
ground up with some gypsum before it is sent to market. 

The carbon dioxide is released into the atmosphere, making it an exothermic reaction. Cement manufacturing
is responsible for about 8\% of the world's $CO_2$ emissions; it is a
major contributor to climate change.

Especially hard concrete, like that used in a nuclear power plant, can
support very large loads without being crushed. However, if
you pull on two ends of a piece of concrete, it comes apart relatively
easily. We say that concrete can handle a lot of \index{compressive stress}\newterm{compressive
  stress}, but not much \index{tensile stress}\newterm{tensile} stress.

\subsection{Steel reinforced concrete}

Many places where we use concrete (like in a bridge), we need both
compressive and tensile stress. Often, the top of a beam is undergoing
compression and the bottom of the beam is undergoing tension.

FIXME Picture here
% https://letstalkscience.ca/educational-resources/backgrounders/why-a-triangle-a-strong-shape

Steel has tremendous tensile strength, but not as much compressive
strength as concrete. To get both tensile \emph{and} compressive
strength, we often bury steel bars or cables inside the concrete.
This is known as \newterm{steel-reinforced concrete}. The concrete
generally does a very good job protecting the steel, which keeps it
from rusting.\index{steel reinforced concrete}

You may have heard of \newterm{rebar} before. That is just short for
``reinforcing bar''.  Typically, rebar has bumps and ridges that keep
the bar and the concrete from moving independently.\index{rebar}

\subsection{Recycling concrete}

Many concrete structures only last about 100 years. When they are
demolished, the concrete can be reused as aggregate in other projects.
Often, the concrete bits are mixed with cement and made into concrete once more.

If the concrete to be reused is reinforced with steel, the steel has
to be removed and recycled separately.  The concrete is then crushed
into small pieces.


\section{Metals}
Elements that transmit electricity well, even at low temperatures, are
called \newterm{metals}\index{metals}. Many metals are likely familiar to you, such as
aluminum, iron, copper, tin, gold, silver, and platinum. Aluminum and
iron are particularly common; together they make up about 14\% of the
earth's crust.

An \newterm{alloy}\index{alloy} is a mixture of elements that includes at least one
metal. Brass, for example, is an alloy of copper and zinc.  Bronze is
an alloy of copper and tin.



\subsection{Steel}

One of the most common alloys is \index{steel}steel, which is an alloy of iron and carbon.
In pure iron, the molecules slip past each other easily, so pure iron
is relatively soft and easily deformed. The carbon in steel prevents
that slipping, which is why steel is much, much harder than iron.

How much carbon does steel have? If you have less than 0.002\% by weight, you end up
with something very much like pure iron.  As you increase the carbon,
it gets harder and harder. Once it gets above about 2\%, the result
is very brittle.

If you add about 11\% chromium to steel, you get \index{stainless steel}\newterm{stainless
  steel}, which resists rusting.

\textbf{Worked example: stress in a steel cable.} A steel cable with a cross section of $2$ mm$^2$ carries a load of $1400$ N. The stress is the force divided by the area:
$$\sigma = \frac{F}{A} = \frac{1400 \text{ N}}{2 \times 10^{-6} \text{ m}^2} = 7.0 \times 10^8 \text{ Pa} = 700 \text{ MPa}$$
This is the tensile strength of the steel listed below, so the cable is at its limit.

\begin{Exercise}[title={Tensile Strength}, label=tensile-mpa]

The \index{tensile strength}tensile strength of steel is usually between 400 MPa and 1200
MPa. A \index{megapascal}Mega Pascal (MPa) is the strength necessary to hold 1,000,000 newtons of
force with a cable that has a 1 square meter cross section. Or,
equivalently, to hold 1 newton of force with a cable that has a 1
square millimeter cross section. 

If you have are buying a round cable that has a tensile strength of
700 Mpa and must hold a 100 kg man aloft, what is the diameter of the
smallest cable you can use?
  
\end{Exercise}
\begin{Answer}[ref=tensile-mpa]
On earth, holding a 100 kg man aloft requires 980 Newtons of force.

$980/700 = 1.4$, so you need a cable with a cross-section area of 1.4
square millimeters.

$$\pi r^2 = 1.4$$

$r = \sqrt{1.4/\pi} \approx .67$ millimeters. This means the cable would
have to have a diameter of at least 1.34 millimeters.

\end{Answer}

\begin{Exercise}[title={Holding Capacity}, label=materials_capacity]
A steel rod with a cross section of $4$ mm$^2$ is loaded to its tensile strength of $700$ MPa. What is the largest force it can carry?
\end{Exercise}
\begin{Answer}[ref=materials_capacity]
$F = \sigma A = (700 \times 10^6)(4 \times 10^{-6}) = 2800$ N.
\end{Answer}

Here are some approximate tensile strengths of other materials:

\begin{tabular}{c|c}
  Material & Tensile strength (MPa) \\
  \hline
  Iron & 3 \\
  Concrete & 4 \\
  Rubber & 16 \\
  Glass & 33 \\
  Wood & 40 \\
  Nylon & 100 \\
  Human hair & 200 \\
  Aluminum  & 300 \\
  Steel & 700 \\
  Spider webs & 1000 \\
  Carbon fiber & 4000
\end{tabular}

\subsection{What metal for what task?}

Copper is often used for electrical wires in your house and
appliances. This is because it is very efficient at moving electricity (very
little power is lost as heat). It is also very good a transmitting
heat, so you will often see copper pots and pans.

Aluminum is less dense than copper, and is still a relatively good
conductor of electricity. This combination of lighter weight and conductivity is why the overhead wires in a power system
are often made of aluminum.

Aluminum is not as strong as steel, but considerably lighter. It is
often used structurally where weight is a concern, such as in skyscrapers, cars,
airplanes, and ships.

\index{titanium}Titanium is about as strong as steel, but it weights about half as
much. Titanium is very difficult to work with, so it is used in places
where weight and strength are very important and cost is not, such as in
airplanes and bicycles.
FIXME: We mention airplanes in both examples. Maybe we should clarify the role each one plays?

(Carbon fiber, which is light, strong, and very easy to work with, is
replacing aluminum and titanium in many applications. 20 years ago,
many expensive bicycles were made of titanium. These days the vast
majority are made with carbon fiber.)

Zinc and tin are very resistant to \index{corrosion}corrosion, so they are often used
as a coating to prevent steel from rusting. They are also used in many
alloys for the same reason.  In the United States, the penny is 97.5\%
zinc and only 2.5\% copper.
